\chapimg{fig/op04.png}
\IfFileExists{extra/lead04.tex}{\chapterlead{\subhd{After this chapter you can} close and audit a water balance; distinguish verification, calibration and validation; compute the bias, mean absolute error, root-mean-square error, coefficient of determination, Nash--Sutcliffe efficiency and index of agreement of a model against data; estimate the uncertainty of a mean and of a regression slope; and state what a one-dimensional model study does and does not claim.}
}{}
\chapter{Auditing a Model: Balance, Statistics and Limits}\label{chap:04}
\markboth{Auditing a Model: Balance, Statistics and Limits}{Auditing a Model: Balance, Statistics and Limits}
\IfFileExists{extra/bg04.tex}{\section{Introduction: why a model needs an audit}
A model is a set of equations that has been put in a computer. Whether its results can be trusted is not a matter of how many equations it has but of what has been checked. An \kw{audit} is a set of such checks. The simplest and the hardest is the closure of the mass balance: what enters the root zone must equal what leaves it plus the change of storage. A model whose balance does not close is losing or creating water, and any saving computed from it is partly an artefact. After closure, the model is compared with data, the comparison is summarised by a few statistics, the uncertainty of the summary is stated, and the limits of what was done are written down. This chapter gives the tools for each of these steps.

\section{Verification, calibration and validation}
Three words are often mixed up. \kw{Verification} asks whether the equations have been solved correctly (is the code right? does it conserve mass? does it reproduce an analytic solution?). \kw{Calibration} adjusts uncertain parameters so that the model fits a set of observations. \kw{Validation} compares the calibrated model with observations \emph{not} used in the calibration. A model that has been verified and calibrated but not validated is a fitted curve; one that has been validated against data from a different site and season is a tested model. In this book the Daisy runs are verified by the balance audit (Chapter~\ref{chap:04}), the application's engines are verified against Daisy (Chapter~9), and no field validation is claimed.

\section{Goodness-of-fit statistics}
Let $O_i$ be $n$ observations (or the reference results) with mean $\bar O$ and $P_i$ the corresponding model values. The commonly used measures are the \kw{bias} $\bar e=\frac1n\sum(P_i-O_i)$, the \kw{mean absolute error} $\mathrm{MAE}=\frac1n\sum|P_i-O_i|$, the \kw{root-mean-square error} $\mathrm{RMSE}=\sqrt{\frac1n\sum(P_i-O_i)^2}$, the \kw{coefficient of determination} $R^2$ of the regression of $P$ on $O$, the \kw{Nash--Sutcliffe efficiency}

\begin{equation}\mathrm{NSE}=1-\frac{\sum_i(O_i-P_i)^2}{\sum_i(O_i-\bar O)^2},\label{eq:nse}\end{equation}
\par\noindent where NSE is the Nash--Sutcliffe efficiency (dimensionless, 1 for a perfect model), $O_i$ the observed and $P_i$ the predicted value at time $i$ and $\bar O$ the mean of the observations.



which is 1 for a perfect model and 0 when the model is no better than the mean of the observations, and the \kw{index of agreement} $d=1-\sum(P_i-O_i)^2/\sum(|P_i-\bar O|+|O_i-\bar O|)^2$ of Willmott, bounded between 0 and 1. No single measure is enough: $R^2$ ignores bias, RMSE is in the units of the data but hides the pattern, and NSE punishes outliers.

\exbox{Statistics of a model against data}{A model predicts the annual evapotranspiration (mm) of eight years: observed 312, 355, 298, 401, 377, 330, 289, 362; predicted 320, 349, 305, 390, 385, 338, 280, 371. Compute the bias, MAE, RMSE, $R^2$, NSE and $d$.}\sol Errors $P-O$: 8, -6, 7, -11, 8, 8, -9, 9. Bias $=1.75$ mm; MAE $=8.25$ mm; RMSE $=\sqrt{70.00}=8.37$ mm. Mean of $O$ is $340.50$; $\sum(O-\bar O)^2=11046$ and $\sum(O-P)^2=560$, so \[\mathrm{NSE}=1-560/11046=0.949\]
Correlation gives $R^2=0.952$; $d=0.987$ \ans The model is nearly unbiased (under 1 mm) and its error is about 2.5 per cent of the mean; NSE of 0.97 says that it explains almost all the variance of the eight years. \finans

\section{Closing the water balance}
For a run of $y$ years the balance of Chapter~2 gives, for each year, the residual $R=P+I-ET-D-Q-\Delta S$ (with $ET=T+E_s+E_o$). A model that conserves mass has $R=0$ to numerical precision. Where $R\ne0$ the audit classifies the cause: a flux that is not reported, a storage that is not counted (snow, ponds, litter, the mulch layer) or a numerical error. The \emph{unattributed evaporative term} of the studies is such a residual that is systematically positive in the mulch experiments: it is evaporation from the mulch layer that the standard output does not count. The balance is the basis of the claim that a saving in one flux is a gain in another, so it is audited first.

\section{Uncertainty of a mean and of a slope}
A mean of $n$ values with standard deviation $s$ has a standard error $s/\sqrt n$ and a 95 per cent confidence interval $\bar x\pm t_{0.975,n-1}s/\sqrt n$. A regression slope $b$ of $y$ on $x$ has standard error $s_b=s_e/\sqrt{\sum(x-\bar x)^2}$ with $s_e$ the residual standard deviation, and the interval $b\pm t\,s_b$. When the distribution is not normal or the statistic is awkward, the \kw{bootstrap} resamples the data with replacement many times and reads the interval from the percentiles of the statistic. An ensemble of simulations has no sampling noise in the usual sense (the same inputs give the same output), so the intervals describe the spread over soils and years, not experimental error; this is stated in every table that quotes them.

\exbox{Slope and its interval}{The data $(x,y)$ are $(1,2.1),(2,3.9),(3,6.2),(4,7.8),(5,10.1),(6,12.2),(7,13.8),(8,16.1)$. Fit a line, find the standard error of the slope and its 95 \% interval ($t_{0.975,6}=2.447$).}\sol \givenline{$t_{0.975,6}=2.447$}\ans $\bar x=4.5$, $\sum(x-\bar x)^2=42$. Slope $b=1.9976$, intercept $a=0.0357$. Residual standard deviation \[s_e=\sqrt{\sum e^2/6}=0.1802\]
\[s_b=0.1802/\sqrt{42}=0.0278\]
The interval is \[1.9976\pm2.447\times0.0278=[1.930,\,2.066]\quad\text{\finans}\]

\section{What a study of this kind can and cannot claim}
The findings of the studies concern a one-dimensional, matrix-flow model with ten soils, a fixed base case of crop and weather, and a designed set of changes. They say what the model does, with the cause separated from the others by the design; they do not say what a field does. Three limits follow. Matrix flow excludes cracks and macropores. One dimension excludes lateral flow and the heterogeneity of a field. A fixed base case limits the extrapolation to other crops and climates. A claim about the \emph{mechanism} (a saving of evaporation is not a saving of water, because the water is conserved) holds for any soil water model that conserves mass; a claim about the \emph{size} of an effect holds only for the model.

}{}
IfFileExists{extra/th04.tex}{\section{Errors, uncertainty and the quality of a measurement}
Every number that comes from a measurement or a model has an error. The \kw{error} is the difference between the value found and the true value, which is never known exactly. Two kinds are distinguished. A \kw{systematic error} (bias) is the same in sign and size every time (a sensor that reads 3 per cent too high); it cannot be reduced by repetition but can be corrected by calibration. A \kw{random error} changes from one reading to the next; it is reduced by averaging. \kw{Accuracy} is the closeness of the result to the true value, and \kw{precision} is the closeness of repeated results to each other: a model can be precise but not accurate.
\par\noindent When a result is computed from several measured values $x_i$ with independent errors $\delta x_i$, the error of $y=f(x_1,x_2,\dots)$ is
\begin{equation*}\boxed{\delta y=\sqrt{\sum_i\left(\frac{\partial f}{\partial x_i}\,\delta x_i\right)^2}}\end{equation*}
\noindent \textbf{Note:} For a sum or difference the absolute errors add in quadrature, and for a product or a quotient the relative errors do. A balance that closes to within the combined error of its terms is satisfactory; a residual larger than that error is a result (Section~\ref{ch:audit-hidden}).

\section{Descriptive statistics}
For a sample $x_1,\dots,x_n$ the \kw{mean} is $\bar x=\sum x_i/n$, the \kw{median} is the middle value, the \kw{variance} is $s^2=\sum(x_i-\bar x)^2/(n-1)$ and the \kw{standard deviation} is $s$. The \kw{percentile} $Q_p$ is the value below which $p$ per cent of the data lie; $Q_{25}$, $Q_{50}$ and $Q_{75}$ are the quartiles, and the interquartile range $\mathrm{IQR}=Q_{75}-Q_{25}$ measures the spread without the effect of extremes. The median and the IQR are preferred when the data are skewed (as annual drainage is, with a few very wet years), and the mean and $s$ when they are symmetric. The \kw{standard error} of the mean, $s/\sqrt n$, measures how well the mean is known, and is not the same as the spread of the data.

\section{Correlation and regression}
The \kw{Pearson correlation coefficient} $r$ measures the strength of a straight-line relation of two variables, from $-1$ to $+1$. The \kw{Spearman rank correlation} $\rho$ is the Pearson coefficient of the ranks; it measures any monotonic relation, is not distorted by outliers and is used in this book (Section~\ref{ch:rank}). \kw{Linear regression} finds the line $y=a+bx$ that minimises the sum of squares of the residuals; the slope is $b=\sum(x_i-\bar x)(y_i-\bar y)/\sum(x_i-\bar x)^2$ and the \kw{coefficient of determination} $R^2$ is the share of the variance of $y$ that the line explains.
\par\noindent \textbf{Notes:} 1. A correlation is not a cause: two quantities can move together because a third drives both. 2. A high $R^2$ for a fitted curve says nothing about the next data set; it must be checked on data left out (cross-validation, Section~\ref{ch:loso}). 3. With ten soils, a correlation of $|\rho|<0.63$ is not significant at the 5 per cent level.

\section{Measures of the agreement of a model with data}
Let $o_i$ be the observed and $m_i$ the modelled values, $i=1,\dots,n$. The usual measures are given below.
\begin{center}\small\begin{tabular}{@{}p{3.8cm}p{5.4cm}p{4.2cm}@{}}\toprule
\hd{Measure} & \hd{Formula} & \hd{Ideal value}\\\midrule
Bias (mean error) & $\dfrac1n\sum(m_i-o_i)$ & 0\\[4pt]
Mean absolute error & $\dfrac1n\sum|m_i-o_i|$ & 0\\[4pt]
Root-mean-square error & $\sqrt{\dfrac1n\sum(m_i-o_i)^2}$ & 0\\[4pt]
Nash--Sutcliffe efficiency & $1-\dfrac{\sum(m_i-o_i)^2}{\sum(o_i-\bar o)^2}$ & 1 ($0$: no better than the mean)\\[4pt]
Index of agreement & $1-\dfrac{\sum(m_i-o_i)^2}{\sum(|m_i-\bar o|+|o_i-\bar o|)^2}$ & 1\\\bottomrule
\end{tabular}\end{center}
\noindent The root-mean-square error is in the unit of the variable and punishes large errors; the Nash--Sutcliffe efficiency is dimensionless, so that variables of different units can be compared. No single measure is enough: a model can have a zero bias and a large RMSE.

\section{Hypothesis tests and the meaning of $p$}
A \kw{hypothesis test} asks whether a difference in the data could be due to chance. The \kw{null hypothesis} states that there is no effect; the \kw{$p$-value} is the probability, if the null hypothesis were true, of a result as extreme as the one observed. A small $p$ (below 0.05) is taken as evidence against the null hypothesis. Two warnings: $p$ is not the probability that the hypothesis is true, and a large sample makes tiny effects ``significant''; the size of the effect (the slope, the difference in per cent) matters more than $p$. A \kw{permutation test}, used in this book, makes the null distribution by shuffling the data many times, and requires no assumption about their distribution.

\section{Confidence intervals and the bootstrap}
A \kw{confidence interval} of 95 per cent is a range that would contain the true value in 95 per cent of repeated experiments. For the mean of a normal sample it is $\bar x\pm t\,s/\sqrt n$. For a statistic with no simple formula the \kw{bootstrap} draws many samples with replacement from the data, computes the statistic for each, and takes the 2.5th and 97.5th percentiles of the results as the interval (Section~\ref{ch:boot}).

\section{Overfitting and the choice of a model}
A model with many parameters can follow the data of the fit and fail on new data; this is \kw{overfitting}. Two safeguards are used: to test on data not used in the fit (leave-one-out, leave-one-soil-out) and to penalise the number of parameters. The \kw{Akaike criterion} $\mathrm{AIC}=n\ln(\mathrm{SSE}/n)+2k$ (for $k$ parameters) is smaller for the better balance between fit and simplicity; when two laws have the same error on new data, the one with fewer parameters is chosen.
}{}
\banner{Closing the water balance: an audit of model outputs}\label{ch:audit}
\section{Why audit a model}
Process-model outputs that support a saving are treated as closed accounts. This can be tested: if the terms of the water balance of Equation~(\ref{eq:bal}) are written for every run, one can check whether they add up, and find out which term is missing when they do not.

\section{A blind search for linear invariants}
Collect, for $N$ runs, the annual terms in a matrix $X$ with one row per run and one column per term, for example $W=P+I$, $ET$, $D$ and $\Delta S+Q$. A \emph{linear invariant} is a vector $c$ with $Xc=0$ for every run. It is found, without being told, as the right singular vector of $X$ with the smallest singular value: if $X=U\Sigma V^{\top}$ then $\|Xc\|$ is smallest for the last column of $V$.

\exbox{Finding the balance blindly}{Four runs have $(W,ET,D,\Delta S+Q)=(800,520,250,30),\ (760,500,240,20),\ (900,560,300,40),\ (700,480,200,20)$ mm. The singular values of $X$ are $1958,\ 43.7,\ 8.01$ and $\approx10^{-13}$. The last is zero to machine precision, and the corresponding singular vector, scaled so that its first entry is 1, is $(1,-1,-1,-1)$: the balance $W=ET+D+(\Delta S+Q)$ was recovered without being supplied. When the study ran this search on the Daisy output, the coefficients were recovered within 5 per cent.}



\begin{kfig}{Closing the balance: the inputs $P+I$ equal the outputs plus the change of storage; what is left after all the reported fluxes is the unattributed term $R$.}
\def\u{0.032}
\draw[axis] (0,0)--(0,4.6);\node[lab,rotate=90] at (-0.7,2.3) {Water (mm)};
\fill[khcyan!35] (0.5,0) rectangle (1.7,{620*\u/3});\fill[khcyan!75] (0.5,{620*\u/3}) rectangle (1.7,{784*\u/3});
\node[lab] at (1.1,{310*\u/3}) {$P$};\node[lab] at (1.1,{702*\u/3}) {$I$};
\fill[khsalmon!60] (3.2,0) rectangle (4.4,{560*\u/3});\fill[khgreen!80!black!40] (3.2,{560*\u/3}) rectangle (4.4,{640*\u/3});\fill[khpurple!70] (3.2,{640*\u/3}) rectangle (4.4,{780*\u/3});\fill[khmag!40] (3.2,{780*\u/3}) rectangle (4.4,{784*\u/3+0.0});
\node[lab] at (3.8,{280*\u/3}) {$ET$};\node[lab] at (3.8,{600*\u/3}) {$D$};\node[lab] at (3.8,{710*\u/3}) {$\Delta S$};\node[lab,right] at (4.5,{784*\u/3}) {$Q,\,R$};
\draw[ln] (0,0)--(5,0);\draw[gd] (0.5,{784*\u/3})--(4.4,{784*\u/3});
\node[lab,below] at (1.1,-0.05) {inputs};\node[lab,below] at (3.8,-0.05) {outputs};
\end{kfig}

\section{The unattributed evaporative term}\label{ch:audit-hidden}
Model outputs report transpiration $T$, soil evaporation $E_s$ and canopy evaporation $E_c$. The non-transpiration evapotranspiration is $ET_{nt}=ET-T$, which in the model is $E_s+E_c$ plus a small residual in the bookkeeping of the evapotranspiration. The audit shows that the residual is not small. Define the unattributed evaporative term
\begin{equation}\label{eq:R}\boxed{R=ET-T-E_s-E_c .}\end{equation}
\par\noindent where $R$ is the residual flux (mm), $ET$ the total evapotranspiration, $T$ the transpiration, $E_s$ the soil evaporation and $E_c$ the canopy evaporation (all in mm over the period).

Averaged over the ensemble, $R=164$ mm\,yr$^{-1}$, which is 30 per cent of the evapotranspiration. It is a property of the model's bookkeeping, not of a real field.

\fig{p1_hidden}{0.95}{The unattributed evaporative term $R$ by experiment (a) and against annual rain (b).}

\section{Judging a saving against the wrong total}
Suppose a mulch or a lower evaporation factor changes the soil evaporation by $\Delta E_s$ and the non-transpiration term by $\Delta ET_{nt}$, and that drainage changes by $\Delta D$. The drainage share of the non-transpiration change is $\psi_D=\Delta D/\Delta ET_{nt}$. If one divides instead by the reported soil evaporation, one obtains $\psi_D'=\Delta D/\Delta E_s=\psi_D\,\Delta ET_{nt}/\Delta E_s$, which differs from $\psi_D$ by the factor $\Delta ET_{nt}/\Delta E_s$.

\exbox{A wrong denominator}{A lower evaporation factor changes soil evaporation by $-40$ mm, and the non-transpiration term by $-55$ mm; drainage changes by $-22$ mm (in magnitudes: drainage given up). The closed share is $\psi_D=22/55=0.40$. Against soil evaporation alone it is $22/40=0.55$, an error of 0.15 --- more than the 0.1 that the study finds in six of ten soils. When $\Delta E_s$ and $\Delta ET_{nt}$ have opposite signs, which happens in two soils, the ratio even changes sign.}



\section{Prediction and its limits}
A compact three-term law predicts evapotranspiration in withheld experiments and withheld soils as well as gradient boosting, with $R^2=0.898$ across withheld experiments and 0.904 across withheld soils. Drainage resists every compact law that was tried. The storage term does not close fully: there is a non-closure of 16 per cent in two experiments.

\begin{daisybox}{reading the logs for the audit}
All terms of this chapter are annual totals read from the Daisy logs: rainfall, irrigation, evapotranspiration, transpiration, soil evaporation, canopy evaporation, drainage at 1 m, storage change and runoff. The matrix $X$ is built from the one-row-per-run summary of the repository. A closed balance means that the singular vector of the smallest singular value is the balance itself.
\end{daisybox}

\exbox{}{Show by hand that $c=(1,-1,-1,-1)$ gives $Xc=0$ for the four runs of the worked example.}\sol \givenline{$c=(1,-1,-1,-1)$ gives ; $Xc=0$}\ans Substituting each row: $800-520-250-30=0$, and so on for the other three rows. \finans

\exbox{}{A run has $ET=540$, $T=230$, $E_s=170$, $E_c=60$. Compute $R$ and the fraction of $ET$ that it represents.}\sol \givenline{$ET=540$ ; $T=230$ ; $E_s=170$ ; $E_c=60$}\ans $R=540-230-170-60=80$ mm, which is $14.8$ per cent of $ET$. \finans

\exbox{}{A saving lowers $E_s$ by 30 mm and $ET_{nt}$ by 36 mm and drainage by 12 mm. Compute $\psi_D$ and $\psi_D'$.}\sol  $\psi_D=12/36=0.33$ and $\psi_D'=12/30=0.40$. \finans

\exbox{}{Why does the success of the blind search on the known balance give confidence in what it reports about $R$?}\sol  Because a method that finds the known balance within 5 per cent is not inventing structure.
\banner{Engines you can check: pipes, design and traceability}
\section{Provenance of equations}
The application has 134 numbered equations. Each carries a label: 38 come from the literature, 18 are data-driven (fitted to the Daisy ensemble) and 78 are engineered for the application. The label is shown beside the number that the equation produced, so a user can find where a number comes from and which origin it has. \finans

\section{Verification}
Every layer has an independent reference test. The design modules reproduce FAO-45 and textbook worked examples within 3 per cent; a meta-model ranks left-out scenarios with Spearman correlations from 0.64 to 0.97 and sweeps $10^{5}$ scenarios in 2 s. Verification found three defects in the first release: a sprinkler range coefficient and a drip iteration count, both corrected, and a fixed advance exponent in the surface layer that under-predicted slow advance by up to 29 per cent (removed in version 1.1.0). It also exposed limits: leaf water potential is set by a conductivity floor, and the drainage share of a mulch saving is understated. A check is valuable only if it can fail, and these failures are reported.

\fig{p2_ws}{0.95}{The surface layer against published advance times and application efficiencies.}

\section{Pipes: Hazen and Williams}
The head loss by friction in a pipe of length $L$, diameter $D$ and flow $Q$ is, in SI units,
\begin{equation}\label{eq:hw}\boxed{h_f=\frac{10.67\,L\,Q^{1.852}}{C^{1.852}D^{4.87}},}\end{equation}
where $C$ is the roughness coefficient of the pipe. The mean speed is $v=4Q/(\pi D^2)$; speed checks keep $v$ in a sensible range. The pump head adds the static lift to the friction losses.

\exbox{A lateral and a pump}{A main line has $L=200$ m, $Q=0.010$ m$^3$\,s$^{-1}$ (10 L\,s$^{-1}$) and $C=140$. With $D=0.10$ m, $h_f=10.67(200)(0.010)^{1.852}/(140^{1.852}\,0.10^{4.87})=3.3$ m and $v=1.27$ m\,s$^{-1}$. With a smaller pipe, $D=0.08$ m, $h_f=9.8$ m and $v=1.99$ m\,s$^{-1}$: the loss is three times larger because $D^{-4.87}$ is steep. The saving of a larger pipe in pump energy has to be weighed against its cost.}



\fig{sketch}{0.95}{The field sketch of SWAFL: seven questions, a drawing area and drawing tools. The pipes drawn here feed the friction and pump-head calculations.}

\exbox{}{Compute $h_f$ for $L=150$ m, $Q=0.006$ m$^3$\,s$^{-1}$, $D=0.075$ m and $C=130$, and the speed $v$.}\sol \givenline{$L=150$ m ; $Q=0.006$ m$^3$\,s$^{-1}$ ; $D=0.075$ m ; $C=130$}\ans $h_f=10.67(150)(0.006)^{1.852}/(130^{1.852}\,0.075^{4.87})$ and $v=4(0.006)/(\pi\,0.075^2)=1.36$ m\,s$^{-1}$; $h_f=4.5$ m. \finans

\exbox{}{By what factor does $h_f$ change if the diameter is doubled? If the flow is doubled?}\sol  Doubling $D$ divides $h_f$ by $2^{4.87}=29.3$; doubling $Q$ multiplies it by $2^{1.852}=3.61$. \finans

\exbox{}{Choose one equation of the application and find its origin label in the application. State whether it is literature, data-driven or engineered, and which experiment of the ensemble would test it.}\sol  Answers depend on the equation chosen.
\banner{Statistics of the evidence}
\section{Reading the statistics of this book in plain words}
Every chapter that follows reports a few numbers that describe how well a law or a model agrees with data: an error, a bias, an $R^2$, a correlation. Each is simple, and none needs advanced mathematics. This section defines them one at a time on a small set of numbers, so that the figures in later chapters can be read at once.

\par\noindent Let there be $N$ pairs of values. $y_i$ is the \emph{observed} value (the number from Daisy or from a measurement) and $\hat y_i$ is the \emph{predicted} value (the number from the law being tested). The \kw{error} of one pair is $e_i=\hat y_i-y_i$: positive when the law over-predicts, negative when it under-predicts.
\begin{center}\small\begin{tabular}{@{}p{2.6cm}p{4.6cm}p{5.6cm}@{}}\toprule
\hd{Measure} & \hd{Formula} & \hd{What it tells you}\\\midrule
Bias & $\bar e=\dfrac1N\sum e_i$ & whether the law is on average too high or too low; zero means the errors cancel\\[4pt]
RMSE & $\sqrt{\dfrac1N\sum e_i^2}$ & the typical size of an error, in the unit of the data; large errors count more because they are squared\\[4pt]
$R^2$ & $1-\dfrac{\sum e_i^2}{\sum(y_i-\bar y)^2}$ & the share of the variation of the data that the law explains; 1 is perfect, 0 is no better than always guessing the mean $\bar y$\\[4pt]
Correlation $r$ & $\dfrac{\sum(y_i-\bar y)(\hat y_i-\bar{\hat y})}{\sqrt{\sum(y_i-\bar y)^2\sum(\hat y_i-\bar{\hat y})^2}}$ & how closely the points follow a straight line; it ignores a constant offset, which bias does not\\\bottomrule
\end{tabular}\end{center}
The reason for two measures of agreement is that they answer different questions. A law can follow every rise and fall of the data ($r$ close to 1) and still be wrong by a constant amount (large bias and RMSE). The RMSE and $R^2$ penalise both faults, and the correlation penalises only the first.

\exbox{Judging a law against six observations}{A law predicts the drained share of irrigation (in mm) for six irrigation events. The observed values are $2.1,\ 3.4,\ 4.0,\ 5.2,\ 6.1,\ 7.5$ and the predicted values are $2.4,\ 3.1,\ 4.5,\ 5.0,\ 6.6,\ 7.0$. Find the bias, the RMSE, $R^2$ and the correlation.}\sol \givenline{$N=6$ ; observed $y=2.1,3.4,4.0,5.2,6.1,7.5$ ; predicted $\hat y=2.4,3.1,4.5,5.0,6.6,7.0$}
\par\noindent\textbf{Step 1: the errors.} $e_i=\hat y_i-y_i$ gives
\[e=+0.3,\ -0.3,\ +0.5,\ -0.2,\ +0.5,\ -0.5\ \text{mm}.\]
\par\noindent\textbf{Step 2: bias.} The sum of the errors is $0.3-0.3+0.5-0.2+0.5-0.5=0.3$, so
\[\bar e=\frac{0.3}{6}=+0.05\ \text{mm}\quad\text{\finans}\]
The law is practically unbiased: over- and under-predictions cancel.
\par\noindent\textbf{Step 3: RMSE.} The squares of the errors are $0.09,\ 0.09,\ 0.25,\ 0.04,\ 0.25,\ 0.25$, with sum $0.97$. Then
\[\mathrm{RMSE}=\sqrt{\frac{0.97}{6}}=\sqrt{0.162}=0.40\ \text{mm}\quad\text{\finans}\]
A typical error is 0.4 mm, against observed values between 2 and 7.5 mm.
\par\noindent\textbf{Step 4: $R^2$.} The mean of the observations is $\bar y=28.3/6=4.72$ mm. The squared distances of the observations from this mean are $6.85,\ 1.73,\ 0.51,\ 0.23,\ 1.91,\ 7.75$, with sum $18.99$. Then
\[R^2=1-\frac{0.97}{18.99}=1-0.051=0.95\quad\text{\finans}\]
The law explains 95 per cent of the variation of the six observations.
\par\noindent\textbf{Step 5: correlation.} Working through the formula in the table gives $r=0.975$, so $r^2=0.951$. Here $r^2$ and $R^2$ are nearly equal because the bias is almost zero. They differ when the law has a large bias, and then $R^2$ is the honest one.

\exbox{The same shape with a constant offset}{Add $+1.5$ mm to every predicted value of the previous example. Find the new bias, RMSE, $R^2$ and correlation, and explain the result.}\sol \givenline{$\hat y_i\to\hat y_i+1.5$}
\par\noindent Every error increases by $1.5$ mm, so the new errors are $1.8,\ 1.2,\ 2.0,\ 1.3,\ 2.0,\ 1.0$. The bias is $\bar e=9.3/6=1.55$ mm. The sum of squares is $3.24+1.44+4.00+1.69+4.00+1.00=15.37$, so $\mathrm{RMSE}=\sqrt{15.37/6}=1.60$ mm and $R^2=1-15.37/18.99=0.19$. The correlation does not change: $r=0.975$, because adding a constant does not change the shape. A law that moves exactly like the data but sits 1.5 mm too high is nearly useless for predicting amounts, and only the bias, the RMSE and $R^2$ show it. \finans

\section{Rank correlation}\label{ch:rank}
Pearson's correlation measures straight-line dependence and is sensitive to outliers. Spearman's rank correlation replaces each value by its rank and is the Pearson correlation of the ranks; when there are no ties
\begin{equation}\label{eq:spearman}\boxed{\rho=1-\frac{6\sum_i d_i^{2}}{N(N^{2}-1)},}\end{equation}
with $d_i$ the difference of the ranks of the $i$-th pair.
\exbox{Memory against $n$ for ten soils}{Rank the ten soils by $n$ (Table~\ref{tab:soils}) and by the memory fate. The rank differences are $+1$ for S01, $+2$ for S02, $-1$ for S06, $-1$ for S08, $-1$ for S09 and zero for the other four, so $\sum d^2=1+4+1+1+1=8$ and $\rho=1-6(8)/(10\cdot99)=0.952$. This is the 0.95 of the study. An exact permutation test, in which the memory values are shuffled over the soils in all $10!$ ways, finds a correlation this high or higher with probability $196/3{,}628{,}800=5.4\times10^{-5}$ (one-sided): only 196 of the permutations have $\sum d^2\le8$.} \finans


\fig{g_mem_n}{0.55}{The memory fate against $n$ for the ten soils: Spearman $\rho=0.95$.}

\section{The bootstrap}\label{ch:boot}
To measure how uncertain a statistic is, resample the data with replacement many times and recompute it. For the mean memory fate of the ten soils, $0.335$, 1,000 resamples (random seed 1) give a 95 per cent interval of $[0.17,\ 0.52]$, the 2.5th and 97.5th percentiles of the resampled means. The study uses the same idea twice: it resamples the 100 lattice cells of each soil, 1,000 times, for the confidence interval of the memory fate, and it resamples soils to find the interval of $\kappa$.

\section{Leave-one-soil-out validation}\label{ch:loso}
A law fitted to all ten soils will always look good on those ten. The honest test is to leave one soil out, fit to the remaining nine, and predict the one that was left out; repeat for each soil and report the error over the predictions.
\exbox{A straight line as a baseline}{Fit a straight line of the memory fate against $n$ ten times, each time without one soil. The prediction errors have a root-mean-square value of $0.12$, against $0.32$ for the benchmark that always predicts the mean of the other nine. A straight line is a crude law; the competition law of Chapter~7, with one parameter, reaches 0.058, so it is more than twice as accurate. (The 0.12 and 0.32 here are our illustration from the numbers of Table~\ref{tab:soils}; 0.058 is the study's value.)}



\section{Fitting a constant to intervals}
For each soil the experiment does not give the reach, only an interval: the reach lies above the largest firing trigger and below the smallest non-firing one, or only above the largest firing trigger if every trigger fires. The constant $\kappa$ is fitted by minimising
\begin{equation}\label{eq:icloss}\boxed{\mathcal L(\kappa)=\sum_{s}\mathrm{dist}\bigl(\log_{10}\hc(\kappa;\alpha_s,n_s),\ [\ell_s,u_s]\bigr)^{2},}\end{equation}
where the distance is zero inside the interval $[\ell_s,u_s]$ (in $\log_{10}$) and is the distance to the nearest end outside it.
\exbox{A toy fit}{Take three soils with synthetic intervals: S03 (sand) $[100,300]$ cm, S05 (sandy loam) $[1000,2000]$ cm and S10 (silty clay loam) bounded only below, $\ge10{,}000$ cm. For $\kappa=10^{-11},\ 10^{-10.5},\ 10^{-10},\ 10^{-9.5},\ 10^{-9},\ 10^{-8.5},\ 10^{-8}$ the loss is $0.079,\ 0.024,\ 0.0013,\ 0,\ 0,\ 0.0003,\ 0.019$. It is exactly zero for $\kappa$ between $1.4\times10^{-10}$ and $2.6\times10^{-9}$: the minimum is a \emph{plateau}, not a point, because all three predictions lie inside their intervals. At $\kappa=10^{-9.5}$ the predicted reaches are $202,\ 1{,}654$ and $132{,}968$ cm. The study's $\kappa=3.47\times10^{-10}$ lies in such a plateau, and its bootstrap interval, $1.0\times10^{-10}$ to $2.8\times10^{-9}$, is of the same width: the data fix $\kappa$ to within about one order of magnitude, and Chapter~5 showed that this is enough, because $\hc\propto\kappa^{-1/E}$.}


\fig{g_kappa}{0.55}{The reach against $\kappa$ for sand, sandy loam and silty clay loam. The dashed line is the fitted value.}

\section{Choosing between laws}
Several laws can be compared on the same footing: fit each to all soils, and for each compute the leave-one-soil-out error and the number of soils whose observed interval contains the prediction. For the four laws of the study the Mualem form has error 0.18 (log$_{10}$) and puts 5 of 10 soils inside; the power law 0.14 and 4; the fixed-saturation law 0.23 and 4; and the absolute-conductivity law is not supported, because $K(\hc)$ would have to span 3.6 decades. The Mualem form is preferred not because it is the smallest error but because it is derived from the physics and has one constant that does not depend on the soil.

\exbox{}{Compute the Spearman correlation of the pairs $(1,2),(2,1),(3,3),(4,5),(5,4)$.}\sol  Rank differences $-1,1,0,-1,1$, $\sum d^2=4$, $\rho=1-24/120=0.8$. \finans

\exbox{}{Explain why the bootstrap resamples cells of the lattice, and not the runs of different lattices.}\sol  Cells within a lattice share a soil and a design, so they are the units that are exchangeable; resampling across different lattices would mix designs. \finans

\exbox{}{Using the toy fit of this chapter, find $\mathcal L$ at $\kappa=3.5\times10^{-10}$ and say whether it is on the plateau.}\sol \givenline{$\kappa=3.5\times10^{-10}$}\ans $\mathcal L(3.5\times10^{-10})=0$; the predicted reaches are $199,\ 1{,}615$ and $127{,}830$ cm, all inside the intervals. It is on the plateau. \finans

\exbox{}{Why is a leave-one-soil-out test harder to pass than a fit to all soils?}\sol  Because the left-out soil cannot influence the fit, so the error measures prediction and not description. \finans

\exbox{}{The exact one-sided permutation $p$ value of the memory--$n$ correlation is $5.4\times10^{-5}$. What is the smallest possible exact one-sided $p$ for $N=10$ with no ties?}\sol \givenline{$N=10$}\ans $1/10!=1/3{,}628{,}800=2.8\times10^{-7}$.
\banner{What is claimed, what is not, and what comes next}
\section{What is shown}
Memory, audit, reach and fate for ten soils in a one-dimensional matrix-flow model. The design modules reproduce published worked examples. Parsers and shared functions agree with independent readers, and defects were removed. Daisy 7.1.14 runs end to end from SWAFL, including a 5.85 ha field from a satellite image to a schedule. Data, code and application are open. \finans

\section{What is not claimed}
\begin{itemize}[leftmargin=*]
\item Field validation: measured drainage, lysimeters and sensor-scheduled plots are still needed. The memory result gives a testable prediction: sand should drain about 85 per cent of the released storage and silty clay loam about 2 per cent.
\item Run-by-run repetition of the published experiments.
\item The Windows build and the online data services on a connected machine.
\item The non-equilibrium criterion $\Pi_\tau$ is proposed, not tested.
\item Several model terms need attention: the baseline offset, sprinkler range coefficients, the wind term and the conductivity floor.
\end{itemize}
The constant $\kappa$ describes the crop, the climate and the sensor, and all E2 runs used one cereal rotation, one weather series and a sensor at 20 cm. For another crop or climate, $\kappa$ is determined again from a trigger experiment on a few soils.

\section{What is covered and what is not yet}
Part~I derives the equations of the manuscripts; Part~II works the mathematics of the SWAFL layers with computed examples: Engine~1 and Engine~2, irrigation control, the field balance and the soil--plant--atmosphere resistances, soil structure and tillage, mulch and evaporation, feedback and the hidden questions, phenotyping, surface, drip and sprinkler design, drainage design, the crop library and crop water requirement, weather and field data (evapotranspiration, satellite indices, runoff and soil loss), the Dashboard, and the governing equations of the Daisy workbench (Richards column, roots, heat, solutes, nitrogen, crop, calibration). Not yet worked in this edition: the reports, workflows and validity chapters of the application; the 107 reference topics, 69 experiment factors and 43 best combinations of the Daisy workbench (tables of Daisy parameters and runs); the screenshot atlas; the full numerical analysis of the Richards solver (only one worked column is shown); nitrogen transformations beyond the forms in Chapter~\ref{ch:daisy}; and a collection of laboratory problems on the open ensemble. The book has not reached the length first planned; it holds what has been worked and checked.
\IfFileExists{extra/app04.tex}{\subsection{The output of the application, panel by panel}
The audit of Chapter 4 is the layer \emph{Full water balance}: every flux of the year, the 20-year means and the closure of the balance are drawn from the output of the model and not from a calculation by hand.

Each column is a year; the stacked bars split the outflow into evaporation, transpiration, interception, runoff and drainage, and the line is the input. The two sides of the balance are equal to within the residual printed below the chart, which is the closure error of this chapter

\appfig{wb2_0}{0.96}{Full water balance, year by year: rain, irrigation, evapotranspiration, drainage and the change of storage for each of the 20 years.}

The bars split the evapotranspiration into transpiration, soil evaporation, interception and the evaporation of the mulch; the sum is checked against the model total.

\appfig{wb2_1}{0.96}{Evapotranspiration split into all its parts (20-year mean): the parts add up to the simulated total.}

The final figure is the mean balance of the 20 years, a single picture of where the water of the field goes and the number to compare with a measured balance.

\appfig{wb2_2}{0.96}{How the water enters and leaves the soil: 20-year means.}

}{}
\IfFileExists{extra/long04.tex}{\subsection{Complete audit of a simulation}
\exbox{Statistics and balance closure of a model run}{A model and a lysimeter gave eight daily evapotranspiration values. Compute the residuals, the bias, the RMSE, the Nash--Sutcliffe efficiency and the correlation, and audit a season's water balance (precipitation 412 mm, irrigation 180 mm, ET 371 mm, drainage 138 mm, storage change 79 mm).}\sol
\ans Data (observed, modelled): (0.82, 0.75), (1.40, 1.55), (2.35, 2.20), (3.10, 3.30), (3.95, 3.80), (4.70, 4.95), (5.20, 5.00), (6.10, 6.30) mm\,d$^{-1}$
\ans Residuals $e_i=M_i-O_i$: -0.07, 0.15, -0.15, 0.20, -0.15, 0.25, -0.20, 0.20
\ans Bias $=\frac{1}{N}\sum e_i=0.029$ mm\,d$^{-1}$
\ans $\sum e_i^2=0.255$, so $RMSE=\sqrt{\sum e_i^2/N}=0.179$ mm\,d$^{-1}$
\ans Mean of the observations $\bar O=3.453$; $\sum(O_i-\bar O)^2=24.349$; Nash--Sutcliffe $NSE=1-\sum e_i^2/\sum(O_i-\bar O)^2=0.9895$
\ans Correlation $r=0.9952$, coefficient of determination $r^2=0.9905$
\ans Relative error $RMSE/\bar O=5.2$ per cent; a model with $NSE>0.9$ and a bias below 5 per cent of the mean is judged good for this purpose
\ans Water-balance audit: inputs $P+I=592$ mm, outputs $ET+D+\Delta S=588$ mm, closure error $4$ mm ($0.7$ per cent of input)


}{}
\IfFileExists{extra/proj04.tex}{\banner{Design project 4: an audit of the water balance of forty runs}
\textbf{Brief.} Forty annual runs of a water-balance model are to be audited before their savings are quoted. For each run the output gives $P$, $I$, $ET$, $D$, $Q$ and the change of storage $\Delta S$ (mm). The task is to compute the residual of every run, summarise it, test it against zero and decide whether the model can be used. the table below shows the first eight runs.

\textbf{Method.} Residual $R_j=P_j+I_j-ET_j-D_j-Q_j-\Delta S_j$; mean $\bar R$, standard deviation $s_R$, and the $t$ statistic $t=\bar R\sqrt N/s_R$ against the hypothesis $\bar R=0$; maximum $|R|$.

\begin{table}[!htbp]\centering\small\caption{First eight of the forty runs of the audit, all in mm; the last column is the residual of the balance.}
\begin{tabular}{lrrrrrrr}\toprule
\hd{Run} & \hd{$P$} & \hd{$I$} & \hd{$ET$} & \hd{$D$} & \hd{$Q$} & \hd{$\Delta S$} & \hd{$R$}\\\midrule
1 & 603 & 0 & 402 & 155 & 17.7 & 28.4 & 0.07\\
2 & 722 & 0 & 511 & 189 & 24.0 & -0.6 & -1.27\\
3 & 710 & 0 & 480 & 165 & 21.6 & 43.9 & -0.19\\
4 & 554 & 0 & 386 & 116 & 17.2 & 33.8 & 1.14\\
5 & 573 & 173 & 519 & 183 & 23.5 & 20.7 & -0.32\\
6 & 553 & 202 & 526 & 185 & 20.5 & 26.1 & -3.96\\
7 & 651 & 0 & 465 & 158 & 18.1 & 12.5 & -1.91\\
8 & 595 & 0 & 399 & 164 & 20.9 & 7.7 & 4.17\\\bottomrule\end{tabular}\end{table}

\textbf{Result.} Over the forty runs $\bar R=-0.345$ mm, $s_R=1.740$ mm, $\max|R|=4.19$ mm and $t=-1.25$; the 95 per cent interval of the mean is $[-0.901,\,0.211]$ mm. The mean residual is not different from zero at the 5 per cent level (the critical $t$ for 39 degrees of freedom is 2.02); the largest single residual is 0.7 per cent of the mean rainfall.

\textbf{Decision.} The residuals are small against the fluxes (a few tenths of a millimetre against hundreds) and have no trend with irrigation; the model conserves water to the precision of its output, and the savings computed from it can be quoted as redistributions of a conserved quantity. Had the residual been systematically positive with the amount of mulch, the same audit would have located an uncounted evaporation, as the unattributed term of this book does.

}{}
\IfFileExists{extra/ten04.tex}{\subsection{Worked solutions on the ten soils}
A season of the bucket is run on each soil and the balance audited, as the application does.

\exbox{Ten-soil example 4a: Coarse loam (S01)}{Coarse loam: run the 125-day maize season (seeded weather, $P=212$ mm, $\Sigma ET_0=797$ mm) on a root zone of 0.6 m ($TAW=90$ mm), irrigating the whole depletion when it reaches 50 per cent of TAW. Audit the balance.}\sol The bucket gives $ET_a=747$ mm, $I=577$ mm in 12 events, drainage $D=45$ mm \ans Balance: $P+I=212+577=789$ mm; $ET+D=747+45=792$ mm \ans Storage change $\Delta S=-3$ mm; residual $=0.0$ mm \ans The closure is satisfied (limit $\pm1$ mm).

\exbox{Ten-soil example 4a: Coarse sand (S02)}{Coarse sand: run the 125-day maize season (seeded weather, $P=212$ mm, $\Sigma ET_0=797$ mm) on a root zone of 0.6 m ($TAW=30$ mm), irrigating the whole depletion when it reaches 50 per cent of TAW. Audit the balance.}\sol The bucket gives $ET_a=747$ mm, $I=595$ mm in 33 events, drainage $D=63$ mm \ans Balance: $P+I=212+595=807$ mm; $ET+D=747+63=810$ mm \ans Storage change $\Delta S=-3$ mm; residual $=0.0$ mm \ans The closure is satisfied (limit $\pm1$ mm).

\exbox{Ten-soil example 4a: Sand (S03)}{Sand: run the 125-day maize season (seeded weather, $P=212$ mm, $\Sigma ET_0=797$ mm) on a root zone of 0.6 m ($TAW=36$ mm), irrigating the whole depletion when it reaches 50 per cent of TAW. Audit the balance.}\sol The bucket gives $ET_a=747$ mm, $I=579$ mm in 26 events, drainage $D=47$ mm \ans Balance: $P+I=212+579=791$ mm; $ET+D=747+47=794$ mm \ans Storage change $\Delta S=-3$ mm; residual $=0.0$ mm \ans The closure is satisfied (limit $\pm1$ mm).

\exbox{Ten-soil example 4a: Loamy sand (S04)}{Loamy sand: run the 125-day maize season (seeded weather, $P=212$ mm, $\Sigma ET_0=797$ mm) on a root zone of 0.6 m ($TAW=54$ mm), irrigating the whole depletion when it reaches 50 per cent of TAW. Audit the balance.}\sol The bucket gives $ET_a=747$ mm, $I=592$ mm in 19 events, drainage $D=59$ mm \ans Balance: $P+I=212+592=803$ mm; $ET+D=747+59=806$ mm \ans Storage change $\Delta S=-3$ mm; residual $=0.0$ mm \ans The closure is satisfied (limit $\pm1$ mm).

\exbox{Ten-soil example 4a: Sandy loam (S05)}{Sandy loam: run the 125-day maize season (seeded weather, $P=212$ mm, $\Sigma ET_0=797$ mm) on a root zone of 0.6 m ($TAW=72$ mm), irrigating the whole depletion when it reaches 50 per cent of TAW. Audit the balance.}\sol The bucket gives $ET_a=747$ mm, $I=557$ mm in 14 events, drainage $D=32$ mm \ans Balance: $P+I=212+557=769$ mm; $ET+D=747+32=779$ mm \ans Storage change $\Delta S=-10$ mm; residual $=0.0$ mm \ans The closure is satisfied (limit $\pm1$ mm).

\exbox{Ten-soil example 4a: Loam (S06)}{Loam: run the 125-day maize season (seeded weather, $P=212$ mm, $\Sigma ET_0=797$ mm) on a root zone of 0.6 m ($TAW=102$ mm), irrigating the whole depletion when it reaches 50 per cent of TAW. Audit the balance.}\sol The bucket gives $ET_a=747$ mm, $I=608$ mm in 11 events, drainage $D=76$ mm \ans Balance: $P+I=212+608=820$ mm; $ET+D=747+76=823$ mm \ans Storage change $\Delta S=-3$ mm; residual $=0.0$ mm \ans The closure is satisfied (limit $\pm1$ mm).

\exbox{Ten-soil example 4a: Silt loam (S07)}{Silt loam: run the 125-day maize season (seeded weather, $P=212$ mm, $\Sigma ET_0=797$ mm) on a root zone of 0.6 m ($TAW=120$ mm), irrigating the whole depletion when it reaches 50 per cent of TAW. Audit the balance.}\sol The bucket gives $ET_a=747$ mm, $I=573$ mm in 9 events, drainage $D=41$ mm \ans Balance: $P+I=212+573=784$ mm; $ET+D=747+41=787$ mm \ans Storage change $\Delta S=-3$ mm; residual $=0.0$ mm \ans The closure is satisfied (limit $\pm1$ mm).

\exbox{Ten-soil example 4a: Sandy clay loam (S08)}{Sandy clay loam: run the 125-day maize season (seeded weather, $P=212$ mm, $\Sigma ET_0=797$ mm) on a root zone of 0.6 m ($TAW=84$ mm), irrigating the whole depletion when it reaches 50 per cent of TAW. Audit the balance.}\sol The bucket gives $ET_a=747$ mm, $I=562$ mm in 12 events, drainage $D=51$ mm \ans Balance: $P+I=212+562=773$ mm; $ET+D=747+51=798$ mm \ans Storage change $\Delta S=-25$ mm; residual $=0.0$ mm \ans The closure is satisfied (limit $\pm1$ mm).

\exbox{Ten-soil example 4a: Clay loam (S09)}{Clay loam: run the 125-day maize season (seeded weather, $P=212$ mm, $\Sigma ET_0=797$ mm) on a root zone of 0.6 m ($TAW=96$ mm), irrigating the whole depletion when it reaches 50 per cent of TAW. Audit the balance.}\sol The bucket gives $ET_a=747$ mm, $I=561$ mm in 11 events, drainage $D=50$ mm \ans Balance: $P+I=212+561=773$ mm; $ET+D=747+50=797$ mm \ans Storage change $\Delta S=-25$ mm; residual $=-0.0$ mm \ans The closure is satisfied (limit $\pm1$ mm).

\exbox{Ten-soil example 4a: Silty clay loam (S10)}{Silty clay loam: run the 125-day maize season (seeded weather, $P=212$ mm, $\Sigma ET_0=797$ mm) on a root zone of 0.6 m ($TAW=102$ mm), irrigating the whole depletion when it reaches 50 per cent of TAW. Audit the balance.}\sol The bucket gives $ET_a=747$ mm, $I=608$ mm in 11 events, drainage $D=76$ mm \ans Balance: $P+I=212+608=820$ mm; $ET+D=747+76=823$ mm \ans Storage change $\Delta S=-3$ mm; residual $=0.0$ mm \ans The closure is satisfied (limit $\pm1$ mm).

}{}
\IfFileExists{extra/tenb04.tex}{\section{Worked solutions on ten crops}
The ten crops of the library differ in their demand, their roots and their sensitivity. The same type of question is answered for each of them at $ET_0=5.5$ mm\,d$^{-1}$; the data are typical values of the crops (Doorenbos--Kassam and FAO 56).

\exbox{Crop example 4a: Maize}{Maize: the seasonal crop evapotranspiration is $\Sigma ET_c\approx0.8K_{c,mid}ET_0L$ for $L=125$ days; rain is 280 mm and the effective share is 0.75. Find the net irrigation requirement and audit the balance with drainage 20 mm and runoff 15 mm.}\sol $\Sigma ET_c=0.8\times1.20\times5.5\times125=660$ mm \ans Effective rain $0.75\times280=210$ mm \ans $NIR=ET_c-P_e=660-210=450$ mm \ans Balance: $P+I=280+450=730$; $ET+D+Q=660+20+15=695$ \ans Residual 35 mm (change of storage) \ans

\exbox{Crop example 4a: Wheat}{Wheat: the seasonal crop evapotranspiration is $\Sigma ET_c\approx0.8K_{c,mid}ET_0L$ for $L=150$ days; rain is 280 mm and the effective share is 0.75. Find the net irrigation requirement and audit the balance with drainage 20 mm and runoff 15 mm.}\sol $\Sigma ET_c=0.8\times1.15\times5.5\times150=759$ mm \ans Effective rain $0.75\times280=210$ mm \ans $NIR=ET_c-P_e=759-210=549$ mm \ans Balance: $P+I=280+549=829$; $ET+D+Q=759+20+15=794$ \ans Residual 35 mm (change of storage) \ans

\exbox{Crop example 4a: Cotton}{Cotton: the seasonal crop evapotranspiration is $\Sigma ET_c\approx0.8K_{c,mid}ET_0L$ for $L=180$ days; rain is 280 mm and the effective share is 0.75. Find the net irrigation requirement and audit the balance with drainage 20 mm and runoff 15 mm.}\sol $\Sigma ET_c=0.8\times1.15\times5.5\times180=911$ mm \ans Effective rain $0.75\times280=210$ mm \ans $NIR=ET_c-P_e=911-210=701$ mm \ans Balance: $P+I=280+701=981$; $ET+D+Q=911+20+15=946$ \ans Residual 35 mm (change of storage) \ans

\exbox{Crop example 4a: Tomato}{Tomato: the seasonal crop evapotranspiration is $\Sigma ET_c\approx0.8K_{c,mid}ET_0L$ for $L=135$ days; rain is 280 mm and the effective share is 0.75. Find the net irrigation requirement and audit the balance with drainage 20 mm and runoff 15 mm.}\sol $\Sigma ET_c=0.8\times1.15\times5.5\times135=683$ mm \ans Effective rain $0.75\times280=210$ mm \ans $NIR=ET_c-P_e=683-210=473$ mm \ans Balance: $P+I=280+473=753$; $ET+D+Q=683+20+15=718$ \ans Residual 35 mm (change of storage) \ans

\exbox{Crop example 4a: Potato}{Potato: the seasonal crop evapotranspiration is $\Sigma ET_c\approx0.8K_{c,mid}ET_0L$ for $L=130$ days; rain is 280 mm and the effective share is 0.75. Find the net irrigation requirement and audit the balance with drainage 20 mm and runoff 15 mm.}\sol $\Sigma ET_c=0.8\times1.15\times5.5\times130=658$ mm \ans Effective rain $0.75\times280=210$ mm \ans $NIR=ET_c-P_e=658-210=448$ mm \ans Balance: $P+I=280+448=728$; $ET+D+Q=658+20+15=693$ \ans Residual 35 mm (change of storage) \ans

\exbox{Crop example 4a: Soybean}{Soybean: the seasonal crop evapotranspiration is $\Sigma ET_c\approx0.8K_{c,mid}ET_0L$ for $L=130$ days; rain is 280 mm and the effective share is 0.75. Find the net irrigation requirement and audit the balance with drainage 20 mm and runoff 15 mm.}\sol $\Sigma ET_c=0.8\times1.15\times5.5\times130=658$ mm \ans Effective rain $0.75\times280=210$ mm \ans $NIR=ET_c-P_e=658-210=448$ mm \ans Balance: $P+I=280+448=728$; $ET+D+Q=658+20+15=693$ \ans Residual 35 mm (change of storage) \ans

\exbox{Crop example 4a: Sugar beet}{Sugar beet: the seasonal crop evapotranspiration is $\Sigma ET_c\approx0.8K_{c,mid}ET_0L$ for $L=180$ days; rain is 280 mm and the effective share is 0.75. Find the net irrigation requirement and audit the balance with drainage 20 mm and runoff 15 mm.}\sol $\Sigma ET_c=0.8\times1.20\times5.5\times180=950$ mm \ans Effective rain $0.75\times280=210$ mm \ans $NIR=ET_c-P_e=950-210=740$ mm \ans Balance: $P+I=280+740=1020$; $ET+D+Q=950+20+15=985$ \ans Residual 35 mm (change of storage) \ans

\exbox{Crop example 4a: Sunflower}{Sunflower: the seasonal crop evapotranspiration is $\Sigma ET_c\approx0.8K_{c,mid}ET_0L$ for $L=125$ days; rain is 280 mm and the effective share is 0.75. Find the net irrigation requirement and audit the balance with drainage 20 mm and runoff 15 mm.}\sol $\Sigma ET_c=0.8\times1.10\times5.5\times125=605$ mm \ans Effective rain $0.75\times280=210$ mm \ans $NIR=ET_c-P_e=605-210=395$ mm \ans Balance: $P+I=280+395=675$; $ET+D+Q=605+20+15=640$ \ans Residual 35 mm (change of storage) \ans

\exbox{Crop example 4a: Onion}{Onion: the seasonal crop evapotranspiration is $\Sigma ET_c\approx0.8K_{c,mid}ET_0L$ for $L=150$ days; rain is 280 mm and the effective share is 0.75. Find the net irrigation requirement and audit the balance with drainage 20 mm and runoff 15 mm.}\sol $\Sigma ET_c=0.8\times1.05\times5.5\times150=693$ mm \ans Effective rain $0.75\times280=210$ mm \ans $NIR=ET_c-P_e=693-210=483$ mm \ans Balance: $P+I=280+483=763$; $ET+D+Q=693+20+15=728$ \ans Residual 35 mm (change of storage) \ans

\exbox{Crop example 4a: Alfalfa}{Alfalfa: the seasonal crop evapotranspiration is $\Sigma ET_c\approx0.8K_{c,mid}ET_0L$ for $L=100$ days; rain is 280 mm and the effective share is 0.75. Find the net irrigation requirement and audit the balance with drainage 20 mm and runoff 15 mm.}\sol $\Sigma ET_c=0.8\times1.15\times5.5\times100=506$ mm \ans Effective rain $0.75\times280=210$ mm \ans $NIR=ET_c-P_e=506-210=296$ mm \ans Balance: $P+I=280+296=576$; $ET+D+Q=506+20+15=541$ \ans Residual 35 mm (change of storage) \ans

}{}
\IfFileExists{extra/teny04.tex}{\subsection{Worked solutions on the ten soils (third scenario)}
The same questions are repeated for different working conditions, to show how the answer moves with them.

\exbox{Ten-soil example 4c: Coarse loam (S01)}{Coarse loam: run the 125-day maize season (seeded weather, $P=147$ mm, $\Sigma ET_0=801$ mm) on a root zone of 0.6 m ($TAW=90$ mm), irrigating the whole depletion when it reaches 70 per cent of TAW. Audit the balance.}\sol The bucket gives $ET_a=711$ mm, $I=582$ mm in 9 events, drainage $D=18$ mm \ans Balance: $P+I=147+582=729$ mm; $ET+D=711+18=730$ mm \ans Storage change $\Delta S=-0$ mm; residual $=-0.0$ mm \ans The closure is satisfied (limit $\pm1$ mm).

\exbox{Ten-soil example 4c: Coarse sand (S02)}{Coarse sand: run the 125-day maize season (seeded weather, $P=147$ mm, $\Sigma ET_0=801$ mm) on a root zone of 0.6 m ($TAW=30$ mm), irrigating the whole depletion when it reaches 70 per cent of TAW. Audit the balance.}\sol The bucket gives $ET_a=716$ mm, $I=592$ mm in 25 events, drainage $D=35$ mm \ans Balance: $P+I=147+592=739$ mm; $ET+D=716+35=751$ mm \ans Storage change $\Delta S=-12$ mm; residual $=0.0$ mm \ans The closure is satisfied (limit $\pm1$ mm).

\exbox{Ten-soil example 4c: Sand (S03)}{Sand: run the 125-day maize season (seeded weather, $P=147$ mm, $\Sigma ET_0=801$ mm) on a root zone of 0.6 m ($TAW=36$ mm), irrigating the whole depletion when it reaches 70 per cent of TAW. Audit the balance.}\sol The bucket gives $ET_a=712$ mm, $I=612$ mm in 22 events, drainage $D=47$ mm \ans Balance: $P+I=147+612=759$ mm; $ET+D=712+47=759$ mm \ans Storage change $\Delta S=-0$ mm; residual $=-0.0$ mm \ans The closure is satisfied (limit $\pm1$ mm).

\exbox{Ten-soil example 4c: Loamy sand (S04)}{Loamy sand: run the 125-day maize season (seeded weather, $P=147$ mm, $\Sigma ET_0=801$ mm) on a root zone of 0.6 m ($TAW=54$ mm), irrigating the whole depletion when it reaches 70 per cent of TAW. Audit the balance.}\sol The bucket gives $ET_a=714$ mm, $I=596$ mm in 15 events, drainage $D=45$ mm \ans Balance: $P+I=147+596=743$ mm; $ET+D=714+45=759$ mm \ans Storage change $\Delta S=-16$ mm; residual $=-0.0$ mm \ans The closure is satisfied (limit $\pm1$ mm).

\exbox{Ten-soil example 4c: Sandy loam (S05)}{Sandy loam: run the 125-day maize season (seeded weather, $P=147$ mm, $\Sigma ET_0=801$ mm) on a root zone of 0.6 m ($TAW=72$ mm), irrigating the whole depletion when it reaches 70 per cent of TAW. Audit the balance.}\sol The bucket gives $ET_a=701$ mm, $I=587$ mm in 11 events, drainage $D=33$ mm \ans Balance: $P+I=147+587=734$ mm; $ET+D=701+33=734$ mm \ans Storage change $\Delta S=-0$ mm; residual $=-0.0$ mm \ans The closure is satisfied (limit $\pm1$ mm).

\exbox{Ten-soil example 4c: Loam (S06)}{Loam: run the 125-day maize season (seeded weather, $P=147$ mm, $\Sigma ET_0=801$ mm) on a root zone of 0.6 m ($TAW=102$ mm), irrigating the whole depletion when it reaches 70 per cent of TAW. Audit the balance.}\sol The bucket gives $ET_a=701$ mm, $I=519$ mm in 7 events, drainage $D=22$ mm \ans Balance: $P+I=147+519=666$ mm; $ET+D=701+22=724$ mm \ans Storage change $\Delta S=-58$ mm; residual $=-0.0$ mm \ans The closure is satisfied (limit $\pm1$ mm).

\exbox{Ten-soil example 4c: Silt loam (S07)}{Silt loam: run the 125-day maize season (seeded weather, $P=147$ mm, $\Sigma ET_0=801$ mm) on a root zone of 0.6 m ($TAW=120$ mm), irrigating the whole depletion when it reaches 70 per cent of TAW. Audit the balance.}\sol The bucket gives $ET_a=705$ mm, $I=520$ mm in 6 events, drainage $D=21$ mm \ans Balance: $P+I=147+520=666$ mm; $ET+D=705+21=726$ mm \ans Storage change $\Delta S=-59$ mm; residual $=0.0$ mm \ans The closure is satisfied (limit $\pm1$ mm).

\exbox{Ten-soil example 4c: Sandy clay loam (S08)}{Sandy clay loam: run the 125-day maize season (seeded weather, $P=147$ mm, $\Sigma ET_0=801$ mm) on a root zone of 0.6 m ($TAW=84$ mm), irrigating the whole depletion when it reaches 70 per cent of TAW. Audit the balance.}\sol The bucket gives $ET_a=705$ mm, $I=552$ mm in 9 events, drainage $D=13$ mm \ans Balance: $P+I=147+552=698$ mm; $ET+D=705+13=718$ mm \ans Storage change $\Delta S=-20$ mm; residual $=0.0$ mm \ans The closure is satisfied (limit $\pm1$ mm).

\exbox{Ten-soil example 4c: Clay loam (S09)}{Clay loam: run the 125-day maize season (seeded weather, $P=147$ mm, $\Sigma ET_0=801$ mm) on a root zone of 0.6 m ($TAW=96$ mm), irrigating the whole depletion when it reaches 70 per cent of TAW. Audit the balance.}\sol The bucket gives $ET_a=705$ mm, $I=554$ mm in 8 events, drainage $D=12$ mm \ans Balance: $P+I=147+554=701$ mm; $ET+D=705+12=717$ mm \ans Storage change $\Delta S=-16$ mm; residual $=0.0$ mm \ans The closure is satisfied (limit $\pm1$ mm).

\exbox{Ten-soil example 4c: Silty clay loam (S10)}{Silty clay loam: run the 125-day maize season (seeded weather, $P=147$ mm, $\Sigma ET_0=801$ mm) on a root zone of 0.6 m ($TAW=102$ mm), irrigating the whole depletion when it reaches 70 per cent of TAW. Audit the balance.}\sol The bucket gives $ET_a=701$ mm, $I=519$ mm in 7 events, drainage $D=22$ mm \ans Balance: $P+I=147+519=666$ mm; $ET+D=701+22=724$ mm \ans Storage change $\Delta S=-58$ mm; residual $=-0.0$ mm \ans The closure is satisfied (limit $\pm1$ mm).

}{}
\IfFileExists{extra/tenby04.tex}{\subsection{Worked solutions on ten crops (hot climate)}
The ten crops of the library differ in their demand, their roots and their sensitivity. The same type of question is answered for each of them at $ET_0=7.5$ mm\,d$^{-1}$; the data are typical values of the crops (Doorenbos--Kassam and FAO 56).

\exbox{Crop example 4c: Maize}{Maize: the seasonal crop evapotranspiration is $\Sigma ET_c\approx0.8K_{c,mid}ET_0L$ for $L=125$ days; rain is 280 mm and the effective share is 0.75. Find the net irrigation requirement and audit the balance with drainage 20 mm and runoff 15 mm.}\sol \givenline{$L=125$ days; rain}\ans $\Sigma ET_c=0.8\times1.20\times7.5\times125=900$ mm \ans Effective rain $0.75\times280=210$ mm \ans $NIR=ET_c-P_e=900-210=690$ mm \ans Balance: $P+I=280+690=970$; $ET+D+Q=900+20+15=935$ \ans Residual 35 mm (change of storage) \finans

\exbox{Crop example 4c: Wheat}{Wheat: the seasonal crop evapotranspiration is $\Sigma ET_c\approx0.8K_{c,mid}ET_0L$ for $L=150$ days; rain is 280 mm and the effective share is 0.75. Find the net irrigation requirement and audit the balance with drainage 20 mm and runoff 15 mm.}\sol \givenline{$L=150$ days; rain}\ans $\Sigma ET_c=0.8\times1.15\times7.5\times150=1035$ mm \ans Effective rain $0.75\times280=210$ mm \ans $NIR=ET_c-P_e=1035-210=825$ mm \ans Balance: $P+I=280+825=1105$; $ET+D+Q=1035+20+15=1070$ \ans Residual 35 mm (change of storage) \finans

\exbox{Crop example 4c: Cotton}{Cotton: the seasonal crop evapotranspiration is $\Sigma ET_c\approx0.8K_{c,mid}ET_0L$ for $L=180$ days; rain is 280 mm and the effective share is 0.75. Find the net irrigation requirement and audit the balance with drainage 20 mm and runoff 15 mm.}\sol \givenline{$L=180$ days; rain}\ans $\Sigma ET_c=0.8\times1.15\times7.5\times180=1242$ mm \ans Effective rain $0.75\times280=210$ mm \ans $NIR=ET_c-P_e=1242-210=1032$ mm \ans Balance: $P+I=280+1032=1312$; $ET+D+Q=1242+20+15=1277$ \ans Residual 35 mm (change of storage) \finans

\exbox{Crop example 4c: Tomato}{Tomato: the seasonal crop evapotranspiration is $\Sigma ET_c\approx0.8K_{c,mid}ET_0L$ for $L=135$ days; rain is 280 mm and the effective share is 0.75. Find the net irrigation requirement and audit the balance with drainage 20 mm and runoff 15 mm.}\sol \givenline{$L=135$ days; rain}\ans $\Sigma ET_c=0.8\times1.15\times7.5\times135=932$ mm \ans Effective rain $0.75\times280=210$ mm \ans $NIR=ET_c-P_e=932-210=722$ mm \ans Balance: $P+I=280+722=1002$; $ET+D+Q=932+20+15=966$ \ans Residual 35 mm (change of storage) \finans

\exbox{Crop example 4c: Potato}{Potato: the seasonal crop evapotranspiration is $\Sigma ET_c\approx0.8K_{c,mid}ET_0L$ for $L=130$ days; rain is 280 mm and the effective share is 0.75. Find the net irrigation requirement and audit the balance with drainage 20 mm and runoff 15 mm.}\sol \givenline{$L=130$ days; rain}\ans $\Sigma ET_c=0.8\times1.15\times7.5\times130=897$ mm \ans Effective rain $0.75\times280=210$ mm \ans $NIR=ET_c-P_e=897-210=687$ mm \ans Balance: $P+I=280+687=967$; $ET+D+Q=897+20+15=932$ \ans Residual 35 mm (change of storage) \finans

\exbox{Crop example 4c: Soybean}{Soybean: the seasonal crop evapotranspiration is $\Sigma ET_c\approx0.8K_{c,mid}ET_0L$ for $L=130$ days; rain is 280 mm and the effective share is 0.75. Find the net irrigation requirement and audit the balance with drainage 20 mm and runoff 15 mm.}\sol \givenline{$L=130$ days; rain}\ans $\Sigma ET_c=0.8\times1.15\times7.5\times130=897$ mm \ans Effective rain $0.75\times280=210$ mm \ans $NIR=ET_c-P_e=897-210=687$ mm \ans Balance: $P+I=280+687=967$; $ET+D+Q=897+20+15=932$ \ans Residual 35 mm (change of storage) \finans

\exbox{Crop example 4c: Sugar beet}{Sugar beet: the seasonal crop evapotranspiration is $\Sigma ET_c\approx0.8K_{c,mid}ET_0L$ for $L=180$ days; rain is 280 mm and the effective share is 0.75. Find the net irrigation requirement and audit the balance with drainage 20 mm and runoff 15 mm.}\sol \givenline{$L=180$ days; rain}\ans $\Sigma ET_c=0.8\times1.20\times7.5\times180=1296$ mm \ans Effective rain $0.75\times280=210$ mm \ans $NIR=ET_c-P_e=1296-210=1086$ mm \ans Balance: $P+I=280+1086=1366$; $ET+D+Q=1296+20+15=1331$ \ans Residual 35 mm (change of storage) \finans

\exbox{Crop example 4c: Sunflower}{Sunflower: the seasonal crop evapotranspiration is $\Sigma ET_c\approx0.8K_{c,mid}ET_0L$ for $L=125$ days; rain is 280 mm and the effective share is 0.75. Find the net irrigation requirement and audit the balance with drainage 20 mm and runoff 15 mm.}\sol \givenline{$L=125$ days; rain}\ans $\Sigma ET_c=0.8\times1.10\times7.5\times125=825$ mm \ans Effective rain $0.75\times280=210$ mm \ans $NIR=ET_c-P_e=825-210=615$ mm \ans Balance: $P+I=280+615=895$; $ET+D+Q=825+20+15=860$ \ans Residual 35 mm (change of storage) \finans

\exbox{Crop example 4c: Onion}{Onion: the seasonal crop evapotranspiration is $\Sigma ET_c\approx0.8K_{c,mid}ET_0L$ for $L=150$ days; rain is 280 mm and the effective share is 0.75. Find the net irrigation requirement and audit the balance with drainage 20 mm and runoff 15 mm.}\sol \givenline{$L=150$ days; rain}\ans $\Sigma ET_c=0.8\times1.05\times7.5\times150=945$ mm \ans Effective rain $0.75\times280=210$ mm \ans $NIR=ET_c-P_e=945-210=735$ mm \ans Balance: $P+I=280+735=1015$; $ET+D+Q=945+20+15=980$ \ans Residual 35 mm (change of storage) \finans

\exbox{Crop example 4c: Alfalfa}{Alfalfa: the seasonal crop evapotranspiration is $\Sigma ET_c\approx0.8K_{c,mid}ET_0L$ for $L=100$ days; rain is 280 mm and the effective share is 0.75. Find the net irrigation requirement and audit the balance with drainage 20 mm and runoff 15 mm.}\sol \givenline{$L=100$ days; rain}\ans $\Sigma ET_c=0.8\times1.15\times7.5\times100=690$ mm \ans Effective rain $0.75\times280=210$ mm \ans $NIR=ET_c-P_e=690-210=480$ mm \ans Balance: $P+I=280+480=760$; $ET+D+Q=690+20+15=725$ \ans Residual 35 mm (change of storage) \finans

}{}
\IfFileExists{extra/tenx04.tex}{\section{Worked solutions on the ten soils (second scenario)}
The same questions are repeated for different working conditions, to show how the answer moves with them.

\exbox{Ten-soil example 4b: Coarse loam (S01)}{Coarse loam: run the 125-day maize season (seeded weather, $P=216$ mm, $\Sigma ET_0=810$ mm) on a root zone of 0.6 m ($TAW=90$ mm), irrigating the whole depletion when it reaches 40 per cent of TAW. Audit the balance.}\sol The bucket gives $ET_a=759$ mm, $I=569$ mm in 14 events, drainage $D=57$ mm \ans Balance: $P+I=216+569=785$ mm; $ET+D=759+57=816$ mm \ans Storage change $\Delta S=-31$ mm; residual $=0.0$ mm \ans The closure is satisfied (limit $\pm1$ mm).

\exbox{Ten-soil example 4b: Coarse sand (S02)}{Coarse sand: run the 125-day maize season (seeded weather, $P=216$ mm, $\Sigma ET_0=810$ mm) on a root zone of 0.6 m ($TAW=30$ mm), irrigating the whole depletion when it reaches 40 per cent of TAW. Audit the balance.}\sol The bucket gives $ET_a=759$ mm, $I=615$ mm in 38 events, drainage $D=71$ mm \ans Balance: $P+I=216+615=830$ mm; $ET+D=759+71=830$ mm \ans Storage change $\Delta S=0$ mm; residual $=0.0$ mm \ans The closure is satisfied (limit $\pm1$ mm).

\exbox{Ten-soil example 4b: Sand (S03)}{Sand: run the 125-day maize season (seeded weather, $P=216$ mm, $\Sigma ET_0=810$ mm) on a root zone of 0.6 m ($TAW=36$ mm), irrigating the whole depletion when it reaches 40 per cent of TAW. Audit the balance.}\sol The bucket gives $ET_a=759$ mm, $I=596$ mm in 34 events, drainage $D=66$ mm \ans Balance: $P+I=216+596=812$ mm; $ET+D=759+66=825$ mm \ans Storage change $\Delta S=-13$ mm; residual $=0.0$ mm \ans The closure is satisfied (limit $\pm1$ mm).

\exbox{Ten-soil example 4b: Loamy sand (S04)}{Loamy sand: run the 125-day maize season (seeded weather, $P=216$ mm, $\Sigma ET_0=810$ mm) on a root zone of 0.6 m ($TAW=54$ mm), irrigating the whole depletion when it reaches 40 per cent of TAW. Audit the balance.}\sol The bucket gives $ET_a=759$ mm, $I=589$ mm in 24 events, drainage $D=59$ mm \ans Balance: $P+I=216+589=805$ mm; $ET+D=759+59=818$ mm \ans Storage change $\Delta S=-13$ mm; residual $=0.0$ mm \ans The closure is satisfied (limit $\pm1$ mm).

\exbox{Ten-soil example 4b: Sandy loam (S05)}{Sandy loam: run the 125-day maize season (seeded weather, $P=216$ mm, $\Sigma ET_0=810$ mm) on a root zone of 0.6 m ($TAW=72$ mm), irrigating the whole depletion when it reaches 40 per cent of TAW. Audit the balance.}\sol The bucket gives $ET_a=759$ mm, $I=599$ mm in 18 events, drainage $D=69$ mm \ans Balance: $P+I=216+599=815$ mm; $ET+D=759+69=828$ mm \ans Storage change $\Delta S=-13$ mm; residual $=0.0$ mm \ans The closure is satisfied (limit $\pm1$ mm).

\exbox{Ten-soil example 4b: Loam (S06)}{Loam: run the 125-day maize season (seeded weather, $P=216$ mm, $\Sigma ET_0=810$ mm) on a root zone of 0.6 m ($TAW=102$ mm), irrigating the whole depletion when it reaches 40 per cent of TAW. Audit the balance.}\sol The bucket gives $ET_a=759$ mm, $I=574$ mm in 13 events, drainage $D=50$ mm \ans Balance: $P+I=216+574=790$ mm; $ET+D=759+50=809$ mm \ans Storage change $\Delta S=-19$ mm; residual $=0.0$ mm \ans The closure is satisfied (limit $\pm1$ mm).

\exbox{Ten-soil example 4b: Silt loam (S07)}{Silt loam: run the 125-day maize season (seeded weather, $P=216$ mm, $\Sigma ET_0=810$ mm) on a root zone of 0.6 m ($TAW=120$ mm), irrigating the whole depletion when it reaches 40 per cent of TAW. Audit the balance.}\sol The bucket gives $ET_a=759$ mm, $I=558$ mm in 11 events, drainage $D=45$ mm \ans Balance: $P+I=216+558=774$ mm; $ET+D=759+45=805$ mm \ans Storage change $\Delta S=-31$ mm; residual $=0.0$ mm \ans The closure is satisfied (limit $\pm1$ mm).

\exbox{Ten-soil example 4b: Sandy clay loam (S08)}{Sandy clay loam: run the 125-day maize season (seeded weather, $P=216$ mm, $\Sigma ET_0=810$ mm) on a root zone of 0.6 m ($TAW=84$ mm), irrigating the whole depletion when it reaches 40 per cent of TAW. Audit the balance.}\sol The bucket gives $ET_a=759$ mm, $I=570$ mm in 16 events, drainage $D=57$ mm \ans Balance: $P+I=216+570=786$ mm; $ET+D=759+57=816$ mm \ans Storage change $\Delta S=-31$ mm; residual $=0.0$ mm \ans The closure is satisfied (limit $\pm1$ mm).

\exbox{Ten-soil example 4b: Clay loam (S09)}{Clay loam: run the 125-day maize season (seeded weather, $P=216$ mm, $\Sigma ET_0=810$ mm) on a root zone of 0.6 m ($TAW=96$ mm), irrigating the whole depletion when it reaches 40 per cent of TAW. Audit the balance.}\sol The bucket gives $ET_a=759$ mm, $I=582$ mm in 14 events, drainage $D=52$ mm \ans Balance: $P+I=216+582=798$ mm; $ET+D=759+52=811$ mm \ans Storage change $\Delta S=-13$ mm; residual $=0.0$ mm \ans The closure is satisfied (limit $\pm1$ mm).

\exbox{Ten-soil example 4b: Silty clay loam (S10)}{Silty clay loam: run the 125-day maize season (seeded weather, $P=216$ mm, $\Sigma ET_0=810$ mm) on a root zone of 0.6 m ($TAW=102$ mm), irrigating the whole depletion when it reaches 40 per cent of TAW. Audit the balance.}\sol The bucket gives $ET_a=759$ mm, $I=574$ mm in 13 events, drainage $D=50$ mm \ans Balance: $P+I=216+574=790$ mm; $ET+D=759+50=809$ mm \ans Storage change $\Delta S=-19$ mm; residual $=0.0$ mm \ans The closure is satisfied (limit $\pm1$ mm).

}{}
\IfFileExists{extra/tenbx04.tex}{\subsection{Worked solutions on ten crops (cool climate)}
The ten crops of the library differ in their demand, their roots and their sensitivity. The same type of question is answered for each of them at $ET_0=4.0$ mm\,d$^{-1}$; the data are typical values of the crops (Doorenbos--Kassam and FAO 56).

\exbox{Crop example 4b: Maize}{Maize: the seasonal crop evapotranspiration is $\Sigma ET_c\approx0.8K_{c,mid}ET_0L$ for $L=125$ days; rain is 280 mm and the effective share is 0.75. Find the net irrigation requirement and audit the balance with drainage 20 mm and runoff 15 mm.}\sol \givenline{$L=125$ days; rain}\ans \[\Sigma ET_c=0.8\times1.20\times4.0\times125=480\ \text{mm}\] \ans Effective rain $0.75\times280=210$ mm \ans $NIR=ET_c-P_e=480-210=270$ mm \ans Balance: $P+I=280+270=550$; $ET+D+Q=480+20+15=515$ \ans Residual 35 mm (change of storage) \finans

\exbox{Crop example 4b: Wheat}{Wheat: the seasonal crop evapotranspiration is $\Sigma ET_c\approx0.8K_{c,mid}ET_0L$ for $L=150$ days; rain is 280 mm and the effective share is 0.75. Find the net irrigation requirement and audit the balance with drainage 20 mm and runoff 15 mm.}\sol \givenline{$L=150$ days; rain}\ans \[\Sigma ET_c=0.8\times1.15\times4.0\times150=552\ \text{mm}\] \ans Effective rain $0.75\times280=210$ mm \ans $NIR=ET_c-P_e=552-210=342$ mm \ans Balance: $P+I=280+342=622$; $ET+D+Q=552+20+15=587$ \ans Residual 35 mm (change of storage) \finans

\exbox{Crop example 4b: Cotton}{Cotton: the seasonal crop evapotranspiration is $\Sigma ET_c\approx0.8K_{c,mid}ET_0L$ for $L=180$ days; rain is 280 mm and the effective share is 0.75. Find the net irrigation requirement and audit the balance with drainage 20 mm and runoff 15 mm.}\sol \givenline{$L=180$ days; rain}\ans \[\Sigma ET_c=0.8\times1.15\times4.0\times180=662\ \text{mm}\] \ans Effective rain $0.75\times280=210$ mm \ans $NIR=ET_c-P_e=662-210=452$ mm \ans Balance: $P+I=280+452=732$; $ET+D+Q=662+20+15=697$ \ans Residual 35 mm (change of storage) \finans

\exbox{Crop example 4b: Tomato}{Tomato: the seasonal crop evapotranspiration is $\Sigma ET_c\approx0.8K_{c,mid}ET_0L$ for $L=135$ days; rain is 280 mm and the effective share is 0.75. Find the net irrigation requirement and audit the balance with drainage 20 mm and runoff 15 mm.}\sol \givenline{$L=135$ days; rain}\ans \[\Sigma ET_c=0.8\times1.15\times4.0\times135=497\ \text{mm}\] \ans Effective rain $0.75\times280=210$ mm \ans $NIR=ET_c-P_e=497-210=287$ mm \ans Balance: $P+I=280+287=567$; $ET+D+Q=497+20+15=532$ \ans Residual 35 mm (change of storage) \finans

\exbox{Crop example 4b: Potato}{Potato: the seasonal crop evapotranspiration is $\Sigma ET_c\approx0.8K_{c,mid}ET_0L$ for $L=130$ days; rain is 280 mm and the effective share is 0.75. Find the net irrigation requirement and audit the balance with drainage 20 mm and runoff 15 mm.}\sol \givenline{$L=130$ days; rain}\ans \[\Sigma ET_c=0.8\times1.15\times4.0\times130=478\ \text{mm}\] \ans Effective rain $0.75\times280=210$ mm \ans $NIR=ET_c-P_e=478-210=268$ mm \ans Balance: $P+I=280+268=548$; $ET+D+Q=478+20+15=513$ \ans Residual 35 mm (change of storage) \finans

\exbox{Crop example 4b: Soybean}{Soybean: the seasonal crop evapotranspiration is $\Sigma ET_c\approx0.8K_{c,mid}ET_0L$ for $L=130$ days; rain is 280 mm and the effective share is 0.75. Find the net irrigation requirement and audit the balance with drainage 20 mm and runoff 15 mm.}\sol \givenline{$L=130$ days; rain}\ans \[\Sigma ET_c=0.8\times1.15\times4.0\times130=478\ \text{mm}\] \ans Effective rain $0.75\times280=210$ mm \ans $NIR=ET_c-P_e=478-210=268$ mm \ans Balance: $P+I=280+268=548$; $ET+D+Q=478+20+15=513$ \ans Residual 35 mm (change of storage) \finans

\exbox{Crop example 4b: Sugar beet}{Sugar beet: the seasonal crop evapotranspiration is $\Sigma ET_c\approx0.8K_{c,mid}ET_0L$ for $L=180$ days; rain is 280 mm and the effective share is 0.75. Find the net irrigation requirement and audit the balance with drainage 20 mm and runoff 15 mm.}\sol \givenline{$L=180$ days; rain}\ans \[\Sigma ET_c=0.8\times1.20\times4.0\times180=691\ \text{mm}\] \ans Effective rain $0.75\times280=210$ mm \ans $NIR=ET_c-P_e=691-210=481$ mm \ans Balance: $P+I=280+481=761$; $ET+D+Q=691+20+15=726$ \ans Residual 35 mm (change of storage) \finans

\exbox{Crop example 4b: Sunflower}{Sunflower: the seasonal crop evapotranspiration is $\Sigma ET_c\approx0.8K_{c,mid}ET_0L$ for $L=125$ days; rain is 280 mm and the effective share is 0.75. Find the net irrigation requirement and audit the balance with drainage 20 mm and runoff 15 mm.}\sol \givenline{$L=125$ days; rain}\ans \[\Sigma ET_c=0.8\times1.10\times4.0\times125=440\ \text{mm}\] \ans Effective rain $0.75\times280=210$ mm \ans $NIR=ET_c-P_e=440-210=230$ mm \ans Balance: $P+I=280+230=510$; $ET+D+Q=440+20+15=475$ \ans Residual 35 mm (change of storage) \finans

\exbox{Crop example 4b: Onion}{Onion: the seasonal crop evapotranspiration is $\Sigma ET_c\approx0.8K_{c,mid}ET_0L$ for $L=150$ days; rain is 280 mm and the effective share is 0.75. Find the net irrigation requirement and audit the balance with drainage 20 mm and runoff 15 mm.}\sol \givenline{$L=150$ days; rain}\ans \[\Sigma ET_c=0.8\times1.05\times4.0\times150=504\ \text{mm}\] \ans Effective rain $0.75\times280=210$ mm \ans $NIR=ET_c-P_e=504-210=294$ mm \ans Balance: $P+I=280+294=574$; $ET+D+Q=504+20+15=539$ \ans Residual 35 mm (change of storage) \finans

\exbox{Crop example 4b: Alfalfa}{Alfalfa: the seasonal crop evapotranspiration is $\Sigma ET_c\approx0.8K_{c,mid}ET_0L$ for $L=100$ days; rain is 280 mm and the effective share is 0.75. Find the net irrigation requirement and audit the balance with drainage 20 mm and runoff 15 mm.}\sol \givenline{$L=100$ days; rain}\ans \[\Sigma ET_c=0.8\times1.15\times4.0\times100=368\ \text{mm}\] \ans Effective rain $0.75\times280=210$ mm \ans $NIR=ET_c-P_e=368-210=158$ mm \ans Balance: $P+I=280+158=438$; $ET+D+Q=368+20+15=403$ \ans Residual 35 mm (change of storage) \finans

}{}
\IfFileExists{extra/ex04.tex}{\backhead{EXERCISES}{khred}
\begin{exlist}
\item Observed: 20, 25, 30, 35; predicted: 22, 24, 33, 34. Find the bias, MAE and RMSE.\exans{0.75; 1.75; 1.94}
\item Compute the NSE of the same data.\exans{0.880}
\item A run has $P=640$, $I=120$, $ET=585$, $D=140$, $Q=8$, $\Delta S=22$ (mm). Find the residual and say whether the balance closes.\exans{$R=5$ mm; no, 5 mm is missing}
\item The mean residual of 25 runs is 0.8 mm with standard deviation 2.0 mm. Is it different from zero at the 5 \% level ($t_{0.975,24}=2.064$)?\exans{t = 2.0; no}
\item Find the 95 \% interval of a mean of 16 values with mean 74 and standard deviation 8 ($t_{0.975,15}=2.131$).\exans{$[69.74,\ 78.26]$}
\item Compute the bias, RMSE and NSE for observed (410, 455, 390, 500, 470, 430, 380, 460) and predicted (402, 470, 385, 487, 479, 421, 395, 449).\exans{-0.88; 11.13; 0.918}
\item If the NSE of a model is 0, what does it mean? If negative?\exans{as good as the mean; worse than the mean}
\item A slope has $b=0.82$ and standard error 0.05 with 10 degrees of freedom ($t=2.228$). Give the 95 \% interval and say whether $b=1$ is inside it.\exans{$[0.709,\ 0.931]$; no}
\item A mulch experiment shows a positive residual of +30 mm in the mulched runs and 0 in the bare runs. Where could the water have gone?\exans{uncounted evaporation from the mulch layer}
\item Why does a bootstrap interval need no assumption of normality?\exans{it uses the empirical distribution of the data by resampling}
\end{exlist}
}{}
\IfFileExists{extra/dyk04.tex}{\dykband{DO YOU KNOW ?}
\begin{exlist}
\item Distinguish verification, calibration and validation.
\item Why is closing the water balance the first test of a model?
\item What does an NSE of 0.97 mean, and what does it not tell you?
\item Why can $R^2$ be high when a model is badly biased?
\item What is a bootstrap interval, and when is it used?
\item Why are the intervals of an ensemble of simulations not experimental errors?
\item What is the unattributed evaporative term, and how was it found?
\item State three limits of a one-dimensional matrix-flow model study.
\end{exlist}
}{}
\IfFileExists{extra/obj04.tex}{\dykband{OBJECTIVE TYPE QUESTIONS}
\begin{exlist}
\item Calibration means\opt{checking the code}{adjusting parameters to fit data}{testing with new data}{closing the balance}
\item An NSE of 0 means that the model is\opt{perfect}{as good as the mean of the data}{useless always}{unbiased}
\item The 95 \% interval of a mean uses\opt{$s/n$}{$t\,s/\sqrt n$}{$s^2$}{$\bar x^2$}
\item A balance that closes has a residual of\opt{one}{zero}{the rainfall}{the storage}
\item RMSE has the units of\opt{the data}{the square of the data}{none}{percent}
\item Validation uses\opt{the calibration data}{independent data}{no data}{simulated data only}
\item The bootstrap resamples\opt{without replacement}{with replacement}{never}{only the mean}
\item Matrix flow excludes\opt{gravity}{cracks and macropores}{evaporation}{roots}
\item The index of agreement $d$ lies between\opt{$-1$ and 1}{0 and 1}{0 and 100}{$-\infty$ and $\infty$}
\item A systematic positive residual in mulch runs indicates\opt{a numerical error only}{an uncounted evaporation}{a perfect model}{excess rain}
\end{exlist}
\objkey{1.~(b)\quad 2.~(b)\quad 3.~(b)\quad 4.~(b)\quad 5.~(a)\quad 6.~(b)\quad 7.~(b)\quad 8.~(b)\quad 9.~(b)\quad 10.~(b)}
}{}
\swaflvideo{04}
